% ~250 words

This paper presented the first systematic empirical analysis of feature
adoption in the LiaScript ecosystem.
Exploiting the GitHub corpus as the only observable proxy for a
tool with no server-side telemetry, we analysed \num{3317} courses
from \num{307} committers using a four-category feature taxonomy.

The central finding is a pronounced adoption gradient.
Presentation features --- particularly narrator/TTS and structural
animations --- are widely used, as is the macro and import system for
reuse.
Interactive affordances show selective adoption, while advanced embedding
features remain nearly invisible in practice despite being core
components of the LiaScript design.

This gradient mirrors a well-documented pattern in software engineering:
industry reports consistently find that roughly \SI{80}{\percent} of
features in software products are rarely or never used~\cite{pendo2019}.
Our three-group comparison adds nuance to this observation by showing
that different author communities exploit \emph{different} feature
subsets --- MINT-the-GAP's standardised configuration achieves
near-universal adoption of reuse and math features but leaves
presentation capabilities unexplored, while Community authors follow a natural
adoption ladder that plateaus well below the full feature space.
The HCI principle of \emph{progressive disclosure}~\cite{Nielsen2006}
--- revealing complexity only as users need it --- maps directly onto
this ladder and should guide future documentation design.
Notably, prolific authors show the same feature profile as occasional
contributors, indicating that experience alone does not drive
exploration --- a finding that underscores the need for \emph{active}
onboarding rather than relying on organic discovery.

These findings translate into three actionable priorities:
progressive documentation following the adoption ladder,
reduced embedding complexity through better tooling,
and shared configurations that expose rather than hide optional features.
AI-assisted course generation, identified as the top priority in a
complementary user survey, may address all three barriers simultaneously.

A companion study will shift the lens from feature adoption to course
\emph{lifecycle}: how courses evolve over time, how collaboration
manifests through commits and forks, and whether GitHub engagement
metrics (stars, forks) serve as valid indicators of educational impact
in a system that, by design, leaves no usage trace.
